% !PRINTED 340
\weberpage{340}

If one then sets in (5)
\[
x = \alpha' \quad \text{and} \quad x = \beta'
\]
and observes that
\[
\beta - \alpha' = \frac{(\beta - \alpha)f(\beta)}{f(\beta) - f(\alpha)}, \quad \beta - \beta' = \frac{f(\beta)}{f'(\beta)},
\]
it follows that
\[
f(\alpha') < 0, \quad f(\beta') > 0.
\]

The root sought is therefore contained between $\alpha'$ and $\beta'$,
and the interval $(\alpha', \beta')$ is smaller than $(\alpha, \beta)$.

If we put $\alpha', \beta'$ in place of $\alpha, \beta$ in these expressions,
we obtain a new, still narrower interval, and so on.

As for the convergence of this procedure, we can assure
ourselves of it in the following way. Following (6), we set:
\weq{7}{
\frac{\beta' - \alpha'}{\beta - \alpha} = \Phi(\alpha, \beta) = \frac{f(\beta)\{f(\alpha) - f(\beta) + (\beta - \alpha)f'(\beta)\}}{(\beta - \alpha)f'(\beta)\{f(\beta) - f(\alpha)\}}.
}
Numerator and denominator are here divisible by $(\beta - \alpha)^2$, and
if we cancel the common factor $(\beta - \alpha)^2$ and put $\xi, \eta$
for $\alpha, \beta$, we obtain an expression of the form
\[
\Phi(\xi, \eta) = \frac{\psi_1(\xi, \eta)}{\psi_2(\xi, \eta)},
\]
in which $\psi_1$ and $\psi_2$ are integral rational functions of the two variables
$\xi, \eta$.

If $x$ is the root lying in the interval $(\alpha, \beta)$, then by our
assumptions about $f(x)$ the two functions $\psi_1$ and $\psi_2$
do not become zero when
\weq{8}{
\alpha \leqq \xi < x, \quad x < \eta \leqq \beta.
}
The function $\Phi(\xi, \eta)$ is less than 1 for $\xi = \alpha, \eta = \beta$, but also for
every other pair of values in the interval $(\alpha, \beta)$ for which $f(\xi)$
is negative and $f(\eta)$ positive, since every such interval
$(\xi, \eta)$ could be taken in place of $(\alpha, \beta)$.

Now since $\Phi(\xi, \eta)$ is a continuous function of the two variables
$\xi, \eta$ so long as these remain in the region (8), the
function $\Phi(\xi, \eta)$ must have a maximum value in this region,
and this must be less than 1, because $\Phi(\xi, \eta)$
remains less than 1 even in the limiting cases $\xi = x, \eta = x$.

A positive proper fraction $\Theta$ can therefore be assigned
such that
\[
\Phi(\xi, \eta) < \Theta
\]
%CONT>>
